\section{案例研究：真实大模型的并行策略}

\subsection{Megatron-LM：系统化的 3D 并行}

\begin{figure}[H]
\centering
\includegraphics[width=0.95\textwidth]{slides-images/slide-055.jpg}
\caption{Megatron-LM 在不同模型规模下的并行配置与 GPU 利用率\protect\footnotemark}
\end{figure}
\footnotetext{Slides 第56页。来自 Narayanan et al., 2021。}

从 1.7B 到 1T 参数的模型，可以清晰地看到经验法则的应用：
\begin{itemize}
    \item 张量并行从 1 增长到 8（封顶于节点内 GPU 数）
    \item 流水线并行在模型超出单节点内存后启用
    \item 数据并行从最大值开始逐步减小
    \item 所有配置都达到了理论峰值 40--52\% 的 FLOPS 利用率
\end{itemize}

\begin{figure}[H]
\centering
\includegraphics[width=0.95\textwidth]{slides-images/slide-056.jpg}
\caption{3D 并行带来的线性扩展效果与调优细节\protect\footnotemark}
\end{figure}
\footnotetext{Slides 第57页。}

\subsection{近期语言模型的并行策略}

\begin{figure}[H]
\centering
\includegraphics[width=0.95\textwidth]{slides-images/slide-057.jpg}
\caption{近期大模型的并行策略选择\protect\footnotemark}
\end{figure}
\footnotetext{Slides 第58页。}

\begin{table}[H]
\centering
\small
\begin{tabular}{lccc}
\toprule
\textbf{模型} & \textbf{数据并行} & \textbf{张量并行} & \textbf{流水线并行} \\
\midrule
OLMo (7B) & FSDP & -- & -- \\
DeepSeek v1 & ZeRO-1 & 有 & 有 \\
DeepSeek v3 & ZeRO-1 & 64路专家并行 & 16路 \\
Qwen & ZeRO-1 & 有 & 有 \\
Qwen-MoE & ZeRO-1 & 专家并行 & 有 \\
\bottomrule
\end{tabular}
\caption{不同模型的并行策略选择，基本遵循 3D 并行经验法则}
\end{table}

\subsection{Llama 3：Meta 的大规模训练实践}

\begin{figure}[H]
\centering
\includegraphics[width=0.95\textwidth]{slides-images/slide-058.jpg}
\caption{Llama 3 的并行策略与训练配置\protect\footnotemark}
\end{figure}
\footnotetext{Slides 第59页。}

Llama 3 的 405B 模型训练使用了标准的 4D 并行组合：
\begin{itemize}
    \item 张量并行 = 8（节点内）
    \item 流水线并行用于跨节点（视阶段调整）
    \item 上下文并行（CP）仅在长上下文训练阶段启用
    \item 数据并行用于剩余 GPU
\end{itemize}

\begin{warningbox}{大规模训练中的硬件故障}
Llama 3 训练报告显示：
\begin{itemize}
    \item 148 次由 GPU 故障导致的训练中断（占总中断的 30\%）
    \item 32 次计划外硬件维护
    \item 更隐蔽的威胁是\textbf{静默数据损坏}——GPU 不报错但输出错误数据，可能毁掉整个训练运行
\end{itemize}
大规模训练不仅需要高效的并行算法，还需要健壮的容错和检查点机制。
\end{warningbox}

\subsection{Gemma 2（TPU 案例）}

\begin{figure}[H]
\centering
\includegraphics[width=0.95\textwidth]{slides-images/slide-059.jpg}
\caption{Gemma 2 的 TPU 训练配置\protect\footnotemark}
\end{figure}
\footnotetext{Slides 第60页。}

Google 的 Gemma 2 基于 TPU 训练，使用 ZeRO Stage 3（类似 FSDP）+ 模型并行 + 数据并行。TPU 的环面网格使模型并行度可以超过 8，避免了流水线并行的复杂性。

\subsection{本章小结}

从 Megatron-LM 到 Llama 3 到 DeepSeek V3，所有大规模训练都遵循类似的经验法则：张量并行用于节点内、流水线或 FSDP 用于跨节点、数据并行扩展总吞吐。具体的并行度配置取决于模型规模、硬件拓扑和批次大小预算。

%% ========== 总结与延伸 ==========

